% Part: second-order-logic
% Chapter: sol-and-sets
% Section: cardinalities

\documentclass[../../../include/open-logic-section]{subfiles}

\begin{document}

% SEGMENT OLP-0339-S01

\olfileid{sol}{set}{crd}

\olsection{கணங்களின் எண்மைகள்}

\begin{explain}
பொருட்களம் முடிவுறுவதா அல்லது முடிவிலாததா, !!{enumerable} ஆ அல்லது
!!{nonenumerable} ஆ என்று வெளிப்படுத்துவதைப் போலவே,
~$\Domain{M}$ இன் ஓர் உட்கணம் முடிவுறுவதா அல்லது முடிவிலாததா,
!!{enumerable} ஆ அல்லது !!{nonenumerable} ஆ என்ற பண்பையும்
வரையறுக்கலாம்.
\end{explain}

\begin{prop}
வாய்பாடு $\fn{Inf}(X) \ident$
% SOURCE CORRECTION TA-SOL-006: make u an injective self-map of X.
\begin{multline*}
\lexists[u][(\lforall[x][(X(x) \lif X(u(x)))] \land {}
  \lforall[x][\lforall[y][((X(x) \land X(y) \land
      \eq[u(x)][u(y)]) \lif \eq[x][y])]] \land {}\\
  \lexists[y][(X(y) \land \lforall[x][(X(x)
      \lif \eq/[y][u(x)])]])]
\end{multline*}
ஆனது மாறி மதிப்பீடு~$s$ ஐப் பொறுத்து நிறைவுறுத்தப்படுவது,
$s(X)$ முடிவிலாதது எனில், எனில் மட்டுமே.
\end{prop}

\begin{prop}
வாய்பாடு $\fn{Count}(X) \ident $
\begin{multline*}
\lexists[z][\lexists[u][(X(z) \land
    \lforall[x][(X(x) \lif X(u(x)))] \land {} \\ \lforall[Y][((Y(z) \land
      \lforall[x][(Y(x) \lif Y(u(x)))]) \lif X = Y])]])
\end{multline*}
ஆனது மாறி மதிப்பீடு~$s$ ஐப் பொறுத்து நிறைவுறுத்தப்படுவது,
$s(X)$ !!{enumerable} எனில், எனில் மட்டுமே.
\end{prop}

!!{nonenumerable} கணங்கள் உள்ளன என்பதும், உண்மையில் முடிவிலா
அளவுகளுக்கு முடிவிலா பல வேறுபட்ட நிலைகள் உள்ளன என்பதும் கான்டரின்
தேற்றத்திலிருந்து நமக்குத் தெரியும். கணக் கோட்பாடு கணங்களின்
அளவுகளுக்கான ஒரு முழு எண்கணிதத்தை உருவாக்கி, கணங்களுக்கு முடிவிலா
எண்மை எண்களை அளிக்கிறது. முடிவுறு கணங்களின் அளவுகளை அளக்கும் எண்மை
எண்களாக இயல் எண்கள் பயன்படுகின்றன. !!{denumerable} கணங்களின் எண்மை
முதல் முடிவிலா எண்மை; அது~$\aleph_0$ (“அலெஃப்-நாட்” அல்லது
“அலெஃப்-பூச்சியம்”) எனப்படுகிறது. அடுத்த முடிவிலா அளவு~$\aleph_1$.
எண்ணத்தக்கதாக இல்லாமல் (அதாவது, அளவு~$\aleph_0$ ஆக இல்லாமல்)
ஒரு கணம் கொண்டிருக்கக்கூடிய மிகச்சிறிய அளவு இதுவே. “$X$ இன் அளவு
$\aleph_0$” என்பதை
$\fn{Aleph}_0(X) \liff \fn{Inf}(X) \land \fn{Count}(X)$ என்று
வரையறுக்கலாம். $X$ இன் எல்லா உட்கணங்களும் முடிவுறுபவையாகவோ
அளவு~$\aleph_0$ உடையவையாகவோ இருந்து, அக்கணம் தானோ
அளவு~$\aleph_0$ உடையதாக இல்லாவிட்டால், எனில் மட்டுமே,
அதன் அளவு~$\aleph_1$. எனவே இதைப் பின்வரும் வாய்பாட்டால்
வெளிப்படுத்தலாம்:
$\fn{Aleph_1}(X) \ident \lforall[Y][(Y \subseteq X \lif
  (\lnot\fn{Inf}(Y) \lor \fn{Aleph}_0(Y)))] \land \lnot
\fn{Aleph}_0(X)$. அளவு $ \aleph_2$ ஆக இருப்பது இதேபோல்
வரையறுக்கப்படுகிறது; இவ்வாறே தொடர்கிறது.

தனிச்சிறப்பு வாய்ந்த ஓர் அளவு, தொடரகத்தின் எண்மை எனப்படும் அளவு.
இது $\Pow{\Nat}$ இன் அளவு; நிகராக~$\Real$ இன் அளவு.
ஒரு கணம் தொடரகத்தின் அளவுடையது என்பதையும் இரண்டாம்தரத் தருக்கத்தில்
வெளிப்படுத்தலாம்; ஆனால் அதற்கு இன்னும் சிறிது வேலை தேவை.

\end{document}
